\section{Listening}

\subsection{Unit 1}
\subsubsection{Conversation}
\begin{enumerate}
	\item Why does Stacey come to Dr. Pearl's office?
	\begin{choices}
		\item To get permission for sick leave.
		\item \textbf{\color{red}To get permission to quit his class.}
		\item To get permission to attend his class.
		\item To get permission to get her essay back.
	\end{choices}
	
	\item What is worrying Stacey about her studies?
	\begin{choices}
		\item She cannot take a film class next semester.
		\item She cannot cover her humanities requirements.
		\item \textbf{\color{red}She feels that the writing course is too challenging.}
		\item She feels that passing an engineering course is difficult.
	\end{choices}
	
	
	\item What does Dr. Pearl suggest Stacey do?
	\begin{choices}
		\item \textbf{\color{red}Sign up for free tutoring in writing.}
		\item Work with him at his office every day.
		\item Go to the University Writing Center each Friday.
		\item Work with him at The Found Librarian everyday.
	\end{choices}

	\item What is Dr. Pearl's attitude toward Stacey?
	\begin{choices}
		\item \textbf{\color{red}Patient.}
		\item Satisfied.
		\item Indifferent.
		\item Disappointed.
	\end{choices}

\end{enumerate}

\subsubsection{Passage}

\begin{enumerate}
	
	\item What can we learn about confidence from the passage?
	\begin{choices}
		\item \textbf{\color{red}It keeps us striving for success.}
		\item It is just an attitude toward life.
		\item It needs justification from other people.
		\item It can help us find our sense of purpose.
	\end{choices}
	
	\item How can we develop confidence, according to the passage?
	\begin{choices}
		\item \textbf{\color{red}By doing the things we fear.}
		\item By learning from our failures.
		\item By knowing more about ourselves.
		\item By avoiding the feeling of insecurity.
	\end{choices}
	
	\item What should we do if our life is not what we long for?
	\begin{choices}
		\item Avoid making hasty decisions.
		\item \textbf{\color{red}Make changes on a daily basis.}
		\item Reassess our goals and strategies.
		\item Be patient with life and see what happens.
	\end{choices}
	
\end{enumerate}


\subsubsection{Lectures 1}

\begin{enumerate}
	\item What did the speaker discover when teaching math in a public school?
	\begin{choices}
		\item Some students could not finish homework assignments.
		\item Seventh-grade math was very difficult for most students.
		\item Every student could work hard to understand the material.
		\item \textbf{\color{red}IQ wasn't the only factor that influenced students' performance.}
	\end{choices}
	
	\item What topic interested the speaker as a psychologist?
	\begin{choices}
		\item Challenges and mental health.
		\item IQ and academic performance.	
		\item \textbf{\color{red}Factors contributing to success.}
		\item Mental health and environment.
	\end{choices}
	
	\item What does the speaker say about grit?
	\begin{choices}
		\item \textbf{\color{red}We have little knowledge about how to build it.}
		\item Grit means living life passionately like it's a sprint.	
		\item It's better to develop grit in students at a younger age.
		\item Students' grit can improve as their learning ability grows.
	\end{choices}

	
\end{enumerate}

\subsubsection{Lectures 2}

\begin{enumerate}
	\item What can we learn about procrastination from the lecture?
	\begin{choices}
		\item It traps us in a cycle of regrets.
		\item \textbf{\color{red}It weakens our motivation and efficiency.}
		\item It limits our overall creativity in the long run.
		\item It affects our emotional health in the long run.
	\end{choices}
	\item What is the benefit of breaking down tasks into smaller steps?
	\begin{choices}
		\item It can reduce the stress of our work.
		\item It can help us set clear and specific goals.
		\item It enables us to prioritize the achievable tasks.
		\item \textbf{\color{red}It allows us to focus on the most important tasks first.}
	\end{choices}
	\item How can we enhance our concentration?
	\begin{choices}
		\item \textbf{\color{red}By surrounding ourselves with visual cues.}
		\item By setting both short-term and long-term goals.
		\item By building a comfortable and cheerful atmosphere.
		\item By playing motivational and inspirational background music.
	\end{choices}
	\item What can we learn about the Pomodoro Technique from the lecture?
	\begin{choices}
		\item \textbf{\color{red}It is an effective technique for time management.}
		\item It can enhance productivity by reinforcing our goals.
		\item It adds a sense of urgency by increasing stress levels.
		\item It is a strategy used to improve our sense of accomplishment.
	\end{choices}

\end{enumerate}


\newpage
\subsection{Unit 2}

\subsubsection{Conversation}

\begin{enumerate}
	\item Why is Mount Rainier so special to the man?
	\begin{choices}
		\item \textbf{\color{red}It was his father's favorite mountain.}
		\item It is a good place for hiking and camping.
		\item He loves the yellow and blue flowers there.
		\item He and his friends had many happymemories there.
	\end{choices}
	\item How often did the man and his father go hiking and camping on Mount Rainier every year?
	\begin{choices}
		\item Once or twice.
		\item Two or three times.
		\item \textbf{\color{red}Three or four times.}
		\item Four or five times.
	\end{choices}
	\item What can we learn about the man's father from the conversation?
	\begin{choices}
		\item He died at the age of 80.
		\item He had a very tough life.
		\item \textbf{\color{red}He died of a heart attack.}
		\item He enjoyed cooking fish.
	\end{choices}
	\item What is the relationship between the two speakers?
	\begin{choices}
		\item Husband and wife.
		\item \textbf{\color{red}Boss and employee.}
		\item Teacher and student.
		\item Boyfriend and girlfriend.
	\end{choices}
\end{enumerate}


\subsubsection{Passage}

\begin{enumerate}
	\item What will happen if we always think we must do something in a certain way?
	\begin{choices}
		\item We will become less flexible.
		\item \textbf{\color{red}We will experience more stress.}
		\item We will become narrow-minded.
		\item We will lose our ability to think critically.
	\end{choices}
	\item How can we make a large project easier, according to the passage?
	\begin{choices}
		\item \textbf{\color{red}By dividing the project into smaller parts.}
		\item By focusing on manageable parts of the project.
		\item By tackling the most difficult part of the project first.
		\item By taking more time to finish the project step by step.
	\end{choices}
	\item What is the benefit of using our beds for sleeping only?
	\begin{choices}
		\item It improves our memory.
		\item It helps us concentrate better.
		\item \textbf{\color{red}It promotes faster and better sleep.}
		\item It helps develop good sleeping habits.
	\end{choices}
	\item What is the passage mainly about?
	\begin{choices}
		\item Tips for enjoying a peaceful life.
		\item Tips for developing positive thinking.
		\item \textbf{\color{red}Tips for reducing stress levels at school.}
		\item Tips for improving academic performance.
	\end{choices}

\end{enumerate}

\subsubsection{Lectures 1}

\begin{enumerate}
	
	\item What is the first step toward managing jealousy?
	\begin{choices}
		\item Cultivating self-esteem.
		\item \textbf{\color{red}Acknowledging its presence.}
		\item Focusing on personal growth
		\item Stopping making comparisons.
		
	\end{choices}
	\item How can we alleviate jealousy in a relationship?
	\begin{choices}
		\item By building trust with our partner.
		\item By figuring out the reasons behind our jealousy.
		\item \textbf{\color{red}By expressing concerns and listening to our partner.}
		\item By appreciating our own unique qualities and strengths.
	\end{choices}
	\item Why has social media turned into a breeding ground for jealousy?
	\begin{choices}
		\item It leaves us with little time to focus on self-growth.
		\item It discourages us from communicating with others.
		\item It leads to a constant need for acceptance and recognition.
		\item \textbf{\color{red}It leads to misguided perceptions and unhealthy comparisons.}
	\end{choices}


\end{enumerate}


\subsubsection{Lectures 2}

\begin{enumerate}

	\item What can we learn about negative emotions from the passage?
	\begin{choices}
		\item They generally fade away slowly.
		\item They form a large part of our lives.
		\item They are mainly associated with failures.
		\item \textbf{\color{red}They often mark the path to great achievements.}
	\end{choices}
	\item What is the result of avoiding or suppressing negative emotions?
	\begin{choices}
		\item It reduces negative emotions temporarily.
		\item \textbf{\color{red}It prolongs the impact of negative emotions.}
		\item It diminishes our awareness of negative emotions.
		\item It helps us better deal with negative emotions later.
	\end{choices}
	\item When is a suitable time to think about the causes of our negative emotions?
	\begin{choices}
		\item \textbf{\color{red}When we are calm enough.}
		\item When we recognize their presence.
		\item When we can get support from professionals.
		\item When we are worrying about unnecessary things.
	\end{choices}
	\item Why does dealing with negative emotions need patience?
	\begin{choices}
		\item Negative emotions are constantly changing.
		\item It takes time for negative emotions to fade away.
		\item Identifying the causes of negative emotions can be difficult.
		\item \textbf{\color{red}It takes time to find an appropriate way to manage negative emotions.}
	\end{choices}
	
\end{enumerate}

\newpage
\subsection{Unit 3}

\subsubsection{Conversation}

\begin{enumerate}

	\item What kind of person does the woman want for a flatmate?
	\begin{choices}
		\item Very considerate and caring.
		\item Somewhat reserved and quiet.
		\item Pretty easy-going and straightforward.
		\item \textbf{\color{red}Sociable but also aware of personal space.}
	\end{choices}
	\item What do we know about the woman's attitude to cleanliness?
	\begin{choices}
		\item \textbf{\color{red}She cares a lot about being neat and tidy.}
		\item She doesn't care whether her flatmate is tidy or not.
		\item She hopes her flatmate will offer to clean up their room.
		\item She thinks she should share the cleaning with her flatmate.
	\end{choices}
	\item Why does the woman think money is a problem when looking for a flatmate?
	\begin{choices}
		\item She cannot afford to live alone.
		\item \textbf{\color{red}Her previous flatmate often delayed payments.}
		\item She needs the rent to be paid as soon as possible.
		\item She thinks reaching a deal regarding rent is difficult.
	\end{choices}
	\item What will the woman do to make sure her flatmate can afford the rent?
	\begin{choices}
		\item Ask them about their work.
		\item Ask them how much they earn.
		\item Ask them to pay the rent each month on time.
		\item \textbf{\color{red}Ask them whether they can pay three months' rent in advance.}
	\end{choices}
	
\end{enumerate}


\subsubsection{Passage}

\begin{enumerate}
	\item How can we build trust and understanding with our neighbors?
	\begin{choices}
		\item By visiting them frequently.
		\item \textbf{\color{red}By getting to know each other.}
		\item By organizing parties together.
		\item By taking family vacations together.
	\end{choices}
	\item What should we do if we are likely to cause problems for our neighbors?
	\begin{choices}
		\item Address problems appropriately when they occur.
		\item Apologize to our neighbors when problems occur.
		\item \textbf{\color{red}Take action beforehand to avoid potential problems.}
		\item Immediately cease activities that may cause problems.
	\end{choices}
	\item How should we react if our neighbors are bothering us?
	\begin{choices}
		\item We should involve the police for help.
		\item We should wait patiently for their explanations.
		\item \textbf{\color{red}We should express our concerns and discuss solutions together.}
		\item We should involve more family members to address the problem.
	\end{choices}

\end{enumerate}

\subsubsection{Lectures 1}

\begin{enumerate}
	\item Who are the most likely loyal customers of a sandwich shop, according to the lecture?
	\begin{choices}
		\item People who love sandwiches.
		\item People who are offered coupons.
		\item \textbf{\color{red}Local residents living down the street.}
		\item Office workers living in the community.
	\end{choices}
	\item Which of the following methods is crucial to building a genuine connection between a business and the local community?
	\begin{choices}
		\item Offering coupons.
		\item Distributing free lunches.
		\item Handing out free samples.
		\item \textbf{\color{red}Sponsoring local organizations.}
	\end{choices}
	\item Which of the following strategies is effective in reinforcing the connection between a business and the local community?
	\begin{choices}
		\item \textbf{\color{red}Creating loyalty programs.}
		\item Offering more discounts to customers.
		\item Choosing several advertising platforms.
		\item Working with the most suitable marketing company.
	\end{choices}

\end{enumerate}


\subsubsection{Lectures 2}

\begin{enumerate}

	\item Why are well-defined boundaries beneficial?
	\begin{choices}
		\item \textbf{\color{red}They help protect privacy and prevent trespassing.}
		\item They separate one's property from their neighbors'.
		\item They prevent the occurrence of neighborhood disputes.
		\item They enhance safety by defining the zones for specific activities.
	\end{choices}
	\item What does the speaker say about boundaries in the past?
	\begin{choices}
		\item They were solid and well-defined.
		\item They were more visible and tangible.
		\item They were not reliable in terms of safety.
		\item \textbf{\color{red}They were not harmful to neighborly communication.}
	\end{choices}
	\item What is the possible consequence of overemphasizing boundaries in today’s world?
	\begin{choices}
		\item There will be fewer casual interactions.
		\item There will be fewer gatherings and parties.
		\item \textbf{\color{red}There will be fewer free exchanges of ideas.}
		\item There will be a misleading sense of security.
	\end{choices}
	\item How should we address boundaries in modern times, according to the lecture?
	\begin{choices}
		\item By making more efforts to prevent invisible boundaries.
		\item \textbf{\color{red}By overcoming the challenges posed by invisible boundaries.}
		\item By striking a balance between invisible and physical boundaries.
		\item By removing all forms of boundaries to promote communication.
	\end{choices}
	
\end{enumerate}

\newpage
\subsection{Unit 4}

\subsubsection{Conversation}

\begin{enumerate}

	\item What is the man's problem?
	\begin{choices}
		\item He has difficulty writing personal stories.
		\item He thinks the creative writing class is boring.
		\item \textbf{\color{red}He doesn't know what to write for his fictional writing assignment.}
		\item He finds it challenging to write three creative stories in two months.
	\end{choices}
	\item How does the woman usually find inspiration for her new paintings?
	\begin{choices}
		\item She goes to the train station.
		\item \textbf{\color{red}She goes for long walks in nature.}
		\item She engages in conversations with artists.
		\item She observes dramatic goodbyes and romantic reunions.
	\end{choices}
	\item What does the man think of the woman's way of finding inspiration?
	\begin{choices}
		\item It's a useless old trick.
		\item It's a pretty good method.
		\item \textbf{\color{red}It might not work for him.}
		\item It only works for artists and writers.
	\end{choices}
	\item What does the woman advise the man to do at the train station?
	\begin{choices}
		\item Listen to others' conversations.
		\item \textbf{\color{red}Sit in the lobby and watch others.}
		\item Talk with strangers to learn their stories.
		\item Share the plot of his story with passers-by.
	\end{choices}

\end{enumerate}

\subsubsection{Passage}

\begin{enumerate}

	\item What do people disagree about regarding scientific research?
	\begin{choices}
		\item Whether research should be conducted for its immediate benefits.
		\item Whether research should continue even without financial support.
		\item Whether the government should distribute funds for research equally.
		\item \textbf{\color{red}Whether the government should support research with less immediate value.}
	\end{choices}
	\item Why should we value research with unknown benefits?
	\begin{choices}
		\item Such research is also aimed at discovering truths.
		\item Much of such research has proven beneficial to society.
		\item Such research can have a long-term impact on our society.
		\item \textbf{\color{red}Failure to support such research hinders the discovery of new knowledge.}
	\end{choices}
	\item What does the speaker say about computer research?
	\begin{choices}
		\item Computer research was given priority even in the face of opposition.
		\item \textbf{\color{red}Computer research was not well-received at first but has proven beneficial.}
		\item People have always been aware of the practical benefits of computer research.
		\item The progress of computer research was slow at first due to insufficient funding.
	\end{choices}

\end{enumerate}

\subsubsection{Lectures 1}

\begin{enumerate}

	\item Why is daydreaming often criticized?
	\begin{choices}
		\item People often get lost in daydreams.
		\item Too many people indulge in daydreaming.
		\item Daydreaming indicates a lack of productivity.
		\item \textbf{\color{red}People are taught that focus is the key to success.}
	\end{choices}
	\item What is the benefit of daydreaming, according to the lecture?
	\begin{choices}
		\item It makes our brains more active.
		\item \textbf{\color{red}It helps us generate original ideas.}
		\item It creates endless possibilities for success.
		\item It helps us understand the people around us.
	\end{choices}
	\item Why is it necessary to exercise control over daydreaming?
	\begin{choices}
		\item Daydreaming can waste time if it occurs too frequently.
		\item Daydreaming may lead to a loss of one's self-awareness.
		\item \textbf{\color{red}Daydreaming can distract one's attention from the current task.}
		\item Daydreaming may foster a constant desire to escape from the present.
	\end{choices}

\end{enumerate}


\subsubsection{Lectures 2}

\begin{enumerate}
	
	\item What can we learn about being artistic from the lecture?
	\begin{choices}
		\item \textbf{\color{red}It involves developing skills and talents.}
		\item It is the ability to think outside the box.
		\item It is about doing something that's unknown to others.
		\item It means creating for specific listeners, viewers, or users.
	\end{choices}
	\item What is creativity, according to the lecture?
	\begin{choices}
		\item \textbf{\color{red}It involves the development of new ideas.}
		\item It involves self-reflection or self-expression.
		\item It is about turning imaginative ideas into reality.
		\item It is about employing artistic skills to create something.
	\end{choices}
	\item What does the speaker say about innovation?
	\begin{choices}
		\item It is more of an imaginative process than a productive one.
		\item It represents a spontaneous departure from the mainstream.
		\item It often leads to outcomes that can be seen in various forms.
		\item \textbf{\color{red}It is about creating something new that can bring value to others.}
		
	\end{choices}
	\item What is the difference between creativity and innovation?
	\begin{choices}
		\item Innovation is a process; creativity is the final product.
		\item \textbf{\color{red}Creativity helps generate ideas; innovation applies them.}
		\item Innovation is linked to science; creativity is linked to art.
		\item Creativity is inward-focused; innovation is outward-focused.
	\end{choices}

\end{enumerate}

\newpage
\subsection{Unit 5}

\subsubsection{Conversation}

\begin{enumerate}

	\item Why has the man packed up?
	\begin{choices}
		\item He is going to visit a friend.
		\item \textbf{\color{red}He is going on a working holiday.}
		\item He is going on a business trip to France.
		\item He has accepted a full-time job on a farm.
	\end{choices}
	\item What is the man interested in?
	\begin{choices}
		\item Doing housework.
		\item \textbf{\color{red}Doing building work.}
		\item Working in the fields.
		\item Working in the garden.
	\end{choices}
	\item What does the man plan to do on weekends while in France?
	\begin{choices}
		\item Stay in the local town.
		\item Do extra work on the farm.
		\item Go to Paris every weekend.
		\item \textbf{\color{red}Visit a few places in France.}
		
	\end{choices}
	\item What does the woman decide to do, according to the conversation?
	\begin{choices}
		\item Find a place that interests her.
		\item Advise the man to give up his plan.
		\item \textbf{\color{red}Join the man for the unique experience.}
		\item Search the Internet for travel information.
	\end{choices}

	
\end{enumerate}


\subsubsection{Passage}

\begin{enumerate}
	
	\item Why did the speaker become dissatisfied with her job at the company?
	\begin{choices}
		\item \textbf{\color{red}She no longer felt passionate.}
		\item There was no salary increase.
		\item She felt there was no flexibility.
		\item Her boss was not understanding.
	\end{choices}
	\item What can help us achieve a sense of fulfillment, according to the speaker?
	\begin{choices}
		\item Starting our own businesses.
		\item Creating something from scratch.
		\item Turning hobbies into professional pursuits.
		\item \textbf{\color{red}Pursuing anything that reflects our passion.}
	\end{choices}
	\item What job does the speaker find most suitable for herself now?
	\begin{choices}
		\item A business owner.
		\item An online teacher.
		\item \textbf{\color{red}A freelance writer.}
		\item An investment advisor.
	\end{choices}
	
\end{enumerate}

\subsubsection{Lectures 1}

\begin{enumerate}

	\item Why did the speaker find it rewarding to interpret for the Spanish author?
	\begin{choices}
		\item He had a crucial role in the conference.
		\item He was praised by the author for his good work.
		\item \textbf{\color{red}He appreciated the audience's reactions to his work.}
		\item He accurately interpreted the author's ideas and works.
	\end{choices}
	\item Why can being an interpreter broaden one's horizons?
	\begin{choices}
		\item Interpreters are on a never-ending journey to different places.
		\item Interpreters play different roles in various events and conferences.
		\item \textbf{\color{red}Interpreters can meet people from diverse backgrounds and cultures.}
		\item Interpreters can acquire world views distinct from those of other people.
	\end{choices}
	\item Why do interpreters have the privilege of witnessing history unfold before their eyes?
	\begin{choices}
		\item They help organize many high-level conferences.
		\item They help make critical decisions in important negotiations.
		\item They often have the best seats at many high-level conferences.
		\item \textbf{\color{red}They enable effective communication during conferences and negotiations.}
	\end{choices}
	\item How does interpreting contribute to one's personal development?
	\begin{choices}
		\item It makes one understand others better.
		\item It helps one manage emotions effectively.
		\item It helps one develop various speaking styles.
		\item \textbf{\color{red}It helps one maintain composure under pressure.}
	\end{choices}
\end{enumerate}


\subsubsection{Lectures 2}

\begin{enumerate}

	\item What can we learn about the "in-between" space from the lecture?
	\begin{choices}
		\item It leads to increased levels of stress.
		\item \textbf{\color{red}It provides time for detachment and recovery.	}
		\item It blurs the boundaries between work and home life.
		\item It encourages greater mental engagement and productivity.
	\end{choices}
	
	\item According to the lecture, what may undermine the positive effects of commuting?
	\begin{choices}
		\item Focusing on the road map.
		\item Getting off at the wrong bus stop.
		\item Taking a longer commute every day.
		\item \textbf{\color{red}Dwelling on the negatives of the workday.}
	\end{choices}
	
	\item Which of the following is mentioned in the lecture as a means of making our commute more positive and fulfilling?
	\begin{choices}
		\item Spending less time on the road.
		\item \textbf{\color{red}Engaging in conversations with friends.}
		\item Watching the changing scenery along the way.
		\item Choosing comfortable means of transportation.
	\end{choices}
	
\end{enumerate}


\newpage
\subsection{Unit 6}

\subsubsection{Conversation}

\begin{enumerate}
	
	\item What is the man's problem?
	\begin{choices}
		\item \textbf{\color{red}He is struggling with a test.}
		\item He hasn't finished reading his textbooks.
		\item He has trouble writing psychology essays.
		\item He thinks there are too many concepts in the textbook.
		
	\end{choices}
	\item Why doesn't the man want to ask Dr. Smith for help?
	\begin{choices}
		\item He is afraid of Dr. Smith.
		\item \textbf{\color{red}He is too shy and nervous.}
		\item He failed Dr. Smith's class once before.
		\item He is often called on by Dr. Smith in class.
		
	\end{choices}
	\item What can we learn about the teaching assistant from the conversation?
	\begin{choices}
		\item She is a great lecturer.
		\item She has office hours on Wednesdays.
		\item \textbf{\color{red}She is good at making complex concepts easy.}
		\item She is known for being willing to help students.
		
	\end{choices}
	\item What will the man probably do, according to the conversation?
	\begin{choices}
		\item Post his difficulties online for suggestions.
		\item Do some research to prepare for the meeting with Dr. Smith.
		\item Talk to some professors who have posted video lectures online.
		\item \textbf{\color{red}Speak with the teaching assistant and watch some online lectures.}
	\end{choices}
	
\end{enumerate}


\subsubsection{Passage}

\begin{enumerate}
	
	\item Why is it important to consider problems from other people's perspectives?
	\begin{choices}
		\item It helps us agree with other people's points.
		\item It helps us be better understood by other people.
		\item It helps us understand the backgrounds of other people.
		\item \textbf{\color{red}It helps us understand the reasons behind other people's points.}
	\end{choices}
	\item What is a major cause of conflicts in relationships, according to the passage?
	\begin{choices}
		\item We don't respect other people's decisions.
		\item We get upset with other people's misbehavior.
		\item \textbf{\color{red}We expect other people to behave in a certain way.}
		\item We blame other people for their mistakes and limitations.
	\end{choices}
	\item How should we solve problems through communication?
	\begin{choices}
		\item By looking for things that are left unsaid.
		\item By bringing up old conflicts and solving them.
		\item By talking frankly and honestly to earn each other's trust.
		\item \textbf{\color{red}By focusing on positive issues and seeking common ground.}
	\end{choices}
	\item What is the passage mainly about?
	\begin{choices}
		\item Different types of conflicts in our relationships.
		\item \textbf{\color{red}Suggestions on dealing with conflicts in our relationships.}
		\item Suggestions on how to avoid conflicts in our relationships.
		\item Reasons behind various types of conflicts in our relationships.
	\end{choices}

	

\end{enumerate}

\subsubsection{Lectures 1}

\begin{enumerate}
	
	\item Why is styrofoam an environmental problem?
	\begin{choices}
		\item It ends up as litter.
		\item It is impossible to degrade.
		\item \textbf{\color{red}It is slow to degrade and hard to recycle.}
		\item Its manufacturing process is harmful to the environment.
	\end{choices}
	\item How did the team feel at the beginning of the experiment?
	\begin{choices}
		\item They were encouraged by their immediate success.
		\item They were doubtful about the method of their experiment.
		\item \textbf{\color{red}They were discouraged by the numerous failures they faced.}
		\item They were relieved as they found the process easier than expected.
		
	\end{choices}
	\item What is the team's solution for the problem of used styrofoam?
	\begin{choices}
		\item Heating it to make it vaporize.
		\item \textbf{\color{red}Creating activated carbon from it.}
		\item Using several methods to recycle it.
		\item Finding a use for it in manufacturing other products.
	\end{choices}
	
\end{enumerate}


\subsubsection{Lectures 2}

\begin{enumerate}
	
	\item What can we learn about analytic thinkers from the lecture?
	\begin{choices}
		\item They focus on the relationships between objects.
		\item They are good at making well-informed choices.
		\item They place a high value on empathy and creativity.
		\item \textbf{\color{red}They focus on individual objects and their characteristics.}
	\end{choices}
	
	\item In which of the following fields does holistic thinking play an important role?
	\begin{choices}
		\item Finance.
		\item Science.
		\item \textbf{\color{red}Psychology.}
		\item Engineering.
	\end{choices}
	
	\item Which of the following is emphasized by ancient Greek philosophers?
	\begin{choices}
		\item \textbf{\color{red}Logical reasoning.}
		\item Harmony and balance.
		\item The acceptance of changes.
		\item The adaptability of thinking patterns.
	\end{choices}
	
	
	\item What can we learn about the development of thinking styles from the lecture?
	\begin{choices}
		\item People are born either analytic or holistic thinkers.	
		\item \textbf{\color{red}People acquire their thinking styles from the environment.}
		\item People's thinking styles can be influenced by personal preferences.
		\item People develop their thinking styles by learning from their experiences.
	\end{choices}

\end{enumerate}
